\documentclass[11pt,a4paper]{article}

\usepackage[utf8]{inputenc}
\usepackage[T1]{fontenc}
\usepackage{amsmath,amssymb,amsthm}
\usepackage{booktabs}
\usepackage{enumitem}
\usepackage{microtype}
\usepackage[margin=2.5cm]{geometry}
\usepackage{hyperref}
\hypersetup{colorlinks=true,linkcolor=blue,citecolor=blue}

\newcommand{\G}{\mathcal{G}}
\newcommand{\T}{\mathcal{T}}
\newcommand{\range}{\mathrm{range}}

\title{Adversarial Review of DCA-Trie:\\Where the Project Is Wrong,\\
Why It Is Wrong, and What Would Fix It}
\author{Bernard Katamanso}
\date{\today}

\begin{document}
\maketitle

\begin{abstract}
This document is a deliberate, adversarial review of the DCA-Trie
project.  Its purpose is not to defend the work but to find where the
claims, the evidence, or the evaluation are wrong.  It surfaces five
concrete problems: (1) the evaluation metric is lenient; (2) the
empirical cost--benefit of filtering is negative; (3) Freebase ID
filtering silently destroys answer reachability; (4) the ontology is
both too coarse and too incomplete for the gates to be discriminating;
and (5) the claims about the fixed v2 are unverified.  For each we
state the evidence, the mechanism, and what would fix it.
\end{abstract}

%======================================================================
\section{Problem 1: The Evaluation Metric Is Lenient}
%======================================================================

The Hits@1 metric uses substring containment rather than exact
matching.  In \texttt{src/utils/qa\_utils.py}:

\begin{verbatim}
def match(s1, s2):
    s1 = normalize(s1)
    s2 = normalize(s2)
    return s2 in s1      # answer must be a SUBSTRING of prediction
\end{verbatim}

The prediction is scored correct if the normalized answer string
appears anywhere inside it.  Consequences:

\begin{itemize}
  \item A prediction of ``Jamaican'' scores a hit against the
    ground truth ``Jamaican Creole English Language.''
  \item Verbose or partially-correct generations are forgiven.
  \item The number does not measure \emph{faithful} reasoning --- the
    exact thing graph-constrained decoding claims to provide.
\end{itemize}

\textbf{Where this is wrong:} the headline numbers (89\%, 88\%) are
upper bounds on faithful reasoning.  The gap between Hits@1 and
exact-match Hits@1 was never measured.

\textbf{Fix:} report exact-match Hits@1 alongside substring Hits@1, and
report a per-path validity score (is every hop an actual KG edge?).

%======================================================================
\section{Problem 2: The Filtering Cost--Benefit Is Negative}
%======================================================================

The central claim of DCA-Trie v1 is that TypeOracle filtering
``reduces paths while preserving accuracy.''  The evidence says
otherwise:

\begin{center}
\begin{tabular}{@{}lccc@{}}
\toprule
  Dataset & Path reduction & Accuracy $\Delta$ & Time $\Delta$ \\
\midrule
  WebQSP (100) & 13.3\% & $-1.0$pp & $-0.13$s/q \\
  CWQ (100)    & 10.4\% & $-4.0$pp & $-0.10$s/q \\
\bottomrule
\end{tabular}
\end{center}

Filtering removes paths but \emph{also removes correct answers}, and the
accuracy loss grows with graph depth ($-1$pp at 2 hops, $-4$pp at 4
hops).  The time saved ($\sim$0.1 s/q) is negligible relative to the
6+ s/q that beam search costs.

\textbf{Where this is wrong:} the value proposition of v1 is
``cheaper for the same accuracy.''  It delivers ``slightly cheaper and
slightly worse,'' with the accuracy penalty concentrated exactly where
the work claims to shine (deep graphs).

\textbf{Fix:} either (a) show filtering helps on a domain where the
ontology is dense and correct, or (b) reframe the contribution as a
\emph{validation} layer (catch wrong predictions) rather than a
\emph{pruning} layer.

%======================================================================
\section{Problem 3: Freebase ID Filtering Destroys Reachability}
%======================================================================

The v2 fix to the ``entity name vs ID'' mismatch was to \emph{filter
Freebase IDs out of the trie}:

\begin{verbatim}
_FB_ID_RE = re.compile(r"^[gm]\.\w+$")
def _get_gated_paths(nx_graph, head_entity, oracle, ...):
    for neighbor in nx_graph.neighbors(head_entity):
        if _is_freebase_id(neighbor):
            continue   # <-- silently drops the path
        ...
\end{verbatim}

This ``fix'' introduces a \emph{recall} bug.  Measured on the 10-sample
demo set:

\begin{center}
\begin{tabular}{@{}lcc@{}}
\toprule
  & Answers reachable & Lost \\
\midrule
  After TypeOracle gates (v1)        & 94\% & 6\% \\
  After v2 Freebase-ID filtering     & 89\% & 11\% \\
\bottomrule
\end{tabular}
\end{center}

Eight of 72 answers (11\%) become \emph{structurally unreachable}
because the only path to them passes through a Freebase ID entity such
as \texttt{m.04kq1r9}.

\textbf{Where this is wrong:} the ``fix'' guarantees the model can never
generate the correct answer for 11\% of questions.  The trie can only
output names, but the KG stores many facts via IDs.  Filtering the IDs
does not make the model smarter --- it removes the evidence.

\textbf{Fix:} \emph{canonicalization}, not filtering.  Map Freebase IDs
to their human-readable names via the KG's \texttt{common.topic}
metadata before building the trie, so the path \emph{name} is in the
trie while the \emph{edge} still points at the ID entity.

%======================================================================
\section{Problem 4: The Ontology Is Both Coarse and Incomplete}
%======================================================================

The TypeOracle's discriminating power is bounded by two facts:

\subsection{Coarseness}
Nine hand-curated type groups and 38 relations.  The ranges are broad:
\texttt{people.person.nationality} has range = 7 location types.  A gate
that admits anything typed ``Location'' prunes almost nothing.

\subsection{Incompleteness}
In the WebQSP subgraphs, roughly 78\% of nodes are Freebase IDs with
\emph{no} recorded types.  Both gates are \emph{conservative}:
unknown entity $\Rightarrow$ allow.  So the gates only fire on the
minority of entities whose types are known, and even then only when
those types conflict with a relation range.

\textbf{Where this is wrong:} the two sources of information
(hand-curated schema, auto-mined schema) both depend on type coverage,
which is sparse.  The gates are permissive not by design choice but
because the ontology cannot discriminate.  The 13\% reduction is the
\emph{ceiling} of what this ontology can prune --- and it comes at a
negative accuracy cost.

\textbf{Fix:} demonstrate the approach on a \emph{controlled} ontology
(substation IEC 61850, medical SNOMED-CT) where types are complete and
ranges are narrow.  The Freebase experiment can only ever show weak
gates.

%======================================================================
\section{Problem 5: The Fixed v2 Claims Are Unverified}
%======================================================================

Commit \texttt{b27570e} claims: ``v2 beam search + Freebase ID
filtering + demo scripts.''  The intent is ``post-fix v2 should be
comparable to v1.''  The evidence:

\begin{itemize}
  \item The only post-fix run is a 3-sample smoke test (33\%).
  \item The 56\%/54.9\% figures predate the context-accumulation fix.
  \item No evaluation of beam-size $>1$, score tracking, or
    diversification has been run.
  \item The claim ``comparable to v1'' is \emph{assumed}, not shown.
\end{itemize}

\textbf{Where this is wrong:} the project is committing to a narrative
(``v2 is fixed'') before the data exists.  The v2 architecture had
already been declared ``conclusively false'' in the earlier
CHAPTER4\_RESULTS notes; the new fixes address symptoms
(beam search, IDs, scores) but not the architectural problem
(tokenization misalignment in per-hop trie construction).

\textbf{Fix:} run the post-fix evaluation \emph{before} continuing to
write about v2.  If the fix does not recover v1-level accuracy, the
correct conclusion is that iterative per-hop decoding is not viable
with this tokenizer, and the thesis should pivot to the v1 + validation
direction.

%======================================================================
\section{What Is Right}
%======================================================================

To be fair, the adversarial review should state what holds up:

\begin{itemize}
  \item The \textbf{baseline reproduction} (89\% vs published 91.6\%) is
    solid.
  \item The \textbf{beam search finding} (+8--9pp over greedy) is a
    genuine, useful contribution.
  \item The \textbf{conceptual framing} --- well-formed chains (DoG)
    are necessary but not sufficient; semantic gates add a real
    constraint --- is intellectually sound and publishable as a
    position/analysis.
  \item The \textbf{domain-ontology vision} (living graphs, controlled
    schemas) is the correct direction and has not been tested yet.
\end{itemize}

%======================================================================
\section{Recommendation}
%======================================================================

The project's strongest current position is not ``filtering improves
KGQA accuracy.''  It is:

\begin{enumerate}
  \item \textbf{Reproduce} the honest numbers: exact-match Hits@1,
    path-validity, and a per-error-class breakdown (type errors vs.
    retrieval errors vs. extraction errors).
  \item \textbf{Reframe} filtering as \emph{validation}: use TypeOracle
    as a post-hoc correctness check (already measured: catches 23.8\%
    of errors at 95.5\% precision), not as a pruning pre-step.
  \item \textbf{Test} the controlled-ontology thesis on a domain graph,
    where the gates can finally be strong.
  \item \textbf{Retire or rigorously re-test} v2; stop claiming fixes
    without post-fix evaluation.
\end{enumerate}

\end{document}
